\documentclass[../../../include/open-logic-section]{subfiles}

\begin{document}

\olfileid{sth}{ordinals}{replacement}
\olsection{جایگزینی}

در~\olref[sth][ordinals][vn]{sec} معرفی اعداد ترتیبی را با این پیشنهاد
توجیه کردیم که می‌توانیم آن‌ها را نوع‌های ترتیب، یعنی نمایندگان معیار
خوش‌ترتیب‌ها، تلقی کنیم. برای کارکردن این پیشنهاد باید اثبات کنیم که
\emph{هر خوش‌ترتیب با یک عدد ترتیبی هم‌ریخت است}. آنگاه می‌توانیم
$\ordtype{A,<}$ را عدد ترتیبی~$\alpha$ تعریف کنیم به‌طوری‌که
$\tuple{A,<} \isomorphic \alpha$.

متأسفانه، نتیجهٔ مطلوب را با استفاده از اصول موضوعی که تاکنون معرفی
کرده‌ایم \emph{نمی‌توانیم} اثبات کنیم. (دلیلش را در
\olref[replacement][strength]{sec} خواهیم دید؛ اما فعلاً نکته این است که
نمی‌توانیم.) به اندیشهٔ تازه‌ای نیاز داریم، و آن چنین است:

\begin{axiom}[طرحوارهٔ جایگزینی]
برای هر فرمول~$\phi(x,y)$، عبارت زیر یک اصل موضوع است:
\begin{quote}
	برای هر~$A$، اگر~$(\forall x \in A)\lexists![y][\phi(x,y)]$، آنگاه
	$\Setabs{y}{(\exists x \in A)\phi(x,y)}$ وجود دارد.
\end{quote}
\end{axiom}
\noindent
این نیز مانند جداسازی یک طرحواره است: برای هریک از بی‌نهایت
$\phi$ متفاوت، یک اصل موضوع به دست می‌دهد. و می‌توان آن را به‌همان اندازه
درست (و معمولاً) چنین نوشت:

\begin{defish}
برای هر فرمول~$\phi(x,y)$ که «$B$» در آن نیامده است، عبارت زیر یک اصل
موضوع است:
\[
\forall A[(\forall x \in A)\lexists![y][\phi(x,y)] \lif \exists B\forall y (y \in B \liff (\exists x \in A)\phi(x,y))]
\]
\end{defish}

بااین‌حال، در نخستین مواجهه این فرمول بسیار پیچیده به نظر می‌رسد. پیامد
سریع زیر از جایگزینی احتمالاً اندیشهٔ شهودی مورد نظر را
\emph{روشن‌تر} بیان می‌کند:\footnote{ترم عبارتی است که (با هر مقداردهی
به متغیرهای آزادش) دقیقاً یک شیء را مشخص می‌کند. برای نمونه،
«$\emptyset$» ترمی است که مجموعهٔ تهی را مشخص می‌کند؛ «$\{x\}$» ترمی است
که تک‌عضویِ~$x$ را مشخص می‌کند ($x$ هرچه باشد)؛ و «$x \cup y$» ترمی است
که اجتماع~$x$ و~$y$ را مشخص می‌کند (آن‌ها هرچه باشند).}

\begin{cor}
برای هر ترم~$\tau(x)$ و هر مجموعهٔ~$A$، مجموعهٔ زیر وجود دارد:
\[
	\Setabs{\tau(x)}{x \in A} = \Setabs{y}{(\exists x \in A)y = \tau(x)}.
\]
\end{cor}

\begin{proof}
چون~$\tau$ یک \emph{ترم} است،
$\forall x \lexists![y][\tau(x)=y]$. به طریق اولی،
$(\forall x \in A)\lexists![y][\tau(x)=y]$. پس بنا بر جایگزینی،
$\Setabs{y}{(\exists x \in A)\tau(x)=y}$ وجود دارد.
\end{proof}
\noindent
این نشان می‌دهد «جایگزینی» نام مناسبی برای اصل موضوع است: با داشتن
مجموعهٔ~$A$ می‌توانید مجموعهٔ تازهٔ~$\Setabs{\tau(x)}{x \in A}$ را با
جایگزین‌کردن هر عضو~$A$ با تصویرش تحت~$\tau$ بسازید. در واقع، با پیروی از
نمادگذاری تصویر یک مجموعه تحت یک تابع، می‌توانیم
$\funimage{\tau}{A}$ را برای~$\Setabs{\tau(x)}{x \in A}$ بنویسیم.

اما نکتهٔ اساسی این است که~$\tau$ یک \emph{ترم} است. لازم نیست
(نامی برای) یک \emph{تابع} به معنای~\olref[sfr][fun][rel]{sec}، یعنی
مجموعه‌ای معین از زوج‌های مرتب، باشد. زیرا اگر~$f$ تابعی (به آن معنا)
باشد، مجموعهٔ~$\funimage{f}{A}=\Setabs{f(x)}{x \in A}$ صرفاً زیرمجموعه‌ای
خاص از~$\ran{f}$ است و وجود آن پیشاپیش تنها با اصول موضوع~$\Zminus$
تضمین می‌شود.\footnote{کافی است
$\Setabs{y \in \bigcup \bigcup f}{(\exists x \in A)y=f(x)}$ را در نظر
بگیرید.} اما جایگزینی، چنان‌که در~\olref[sth][replacement][]{chap}
خواهیم دید، افزوده‌ای \emph{نیرومند} به اصول موضوع ماست.

\end{document}
